\chapter{Scanning}

\section{The scanning problem}

A \textbf{scanner} (lexer, tokenizer) is the first phase of a compiler. It reads a string of
characters and splits it into \textbf{tokens}, each a pair (\emph{kind}, \emph{lexeme}):
\[
\mathtt{int\ x = 42;} \;\Longrightarrow\;
(\mathrm{INT},\mathtt{int})\ (\mathrm{ID},\mathtt{x})\ (\mathrm{BECOMES},\mathtt{=})\
(\mathrm{NUM},\mathtt{42})\ (\mathrm{SEMI},\mathtt{;})
\]
Whitespace and comments are discarded. Each token kind is a regular language, so the
language of all valid token sequences, $L = (T_1 \mid T_2 \mid \dots \mid T_k)^*$ (with
whitespace interleaved), is regular too.

\begin{important}[Scanning is a decomposition problem]
Recognising whether $w \in L^*$ is easy, but the scanner must also produce \emph{a
decomposition} of $w$ into tokens. In general the decomposition is not unique: is
\code{abc} one identifier or three? Is \code{>=} one token or \code{>} then \code{=}? We need
a rule to choose.
\end{important}

\section{Building a token-recognising DFA}

\begin{enumerate}
  \item Write a regular expression for each token kind.
  \item Build an $\varepsilon$-NFA for their union, remembering for each accepting state which
        token kind it accepts.
  \item Convert to a DFA by subset construction. An accepting DFA state may contain accepting
        NFA states for several kinds (e.g.\ keyword \code{int} vs.\ identifier); resolve by a
        priority order (keywords beat identifiers).
\end{enumerate}

In practice (and in the CS 241 WLP4 scanner assignment) keywords are recognised as
identifiers by the DFA, and afterwards a lookup checks whether the lexeme is a reserved word.

\begin{example}[Partial WLP4 token DFA]
A fragment of the DFA (keywords are recognised as \code{ID} and looked up afterwards):
\begin{center}
\begin{tikzpicture}[shorten >=1pt,node distance=2.4cm,on grid,auto,>=Stealth,every state/.style={minimum size=1cm,font=\scriptsize}]
  \node[state,initial] (s) at (0,0) {$s$};
  \node[state,accepting] (id) at (3.2,1.8) {ID};
  \node[state,accepting] (num) at (3.2,0) {NUM};
  \node[state,accepting] (zero) at (3.2,-1.8) {ZERO};
  \node[state,accepting] (lt) at (0,-3.4) {LT};
  \node[state,accepting] (le) at (3.2,-3.4) {LE};
  \path[->]
    (s) edge node {a-z,A-Z} (id)
    (id) edge[loop right] node {a-z,A-Z,0-9} ()
    (s) edge node {1-9} (num)
    (num) edge[loop right] node {0-9} ()
    (s) edge node[pos=0.7] {0} (zero)
    (s) edge node[swap] {$<$} (lt)
    (lt) edge node {$=$} (le);
\end{tikzpicture}
\end{center}
Note \code{ZERO} has no outgoing edges: in WLP4 \code{0} followed by a digit is an error,
not a longer number.
\end{example}

\section{Maximal munch}

\begin{definition}
The \textbf{maximal munch} rule: at each step, choose the \emph{longest} prefix of the
remaining input that is a valid token.
\end{definition}

\paragraph{General maximal munch algorithm.} Run the DFA from the current position as far as
possible, remembering the last position at which it was in an accepting state. When the DFA
gets stuck (or the input ends), back up to that last accepting position, emit that token,
and restart from there.

\begin{lstlisting}[style=cpp]
std::vector<Token> scan(const std::string &w) {
    std::vector<Token> out;
    size_t pos = 0;
    while (pos < w.size()) {
        int state = START;
        size_t i = pos, lastAcceptEnd = std::string::npos;
        int lastAcceptState = ERROR;
        while (i < w.size() && delta(state, w[i]) != ERROR) {
            state = delta(state, w[i++]);
            if (accepting(state)) { lastAcceptEnd = i; lastAcceptState = state; }
        }
        if (lastAcceptEnd == std::string::npos) throw ScanError(pos);
        out.push_back({kindOf(lastAcceptState), w.substr(pos, lastAcceptEnd - pos)});
        pos = lastAcceptEnd;          // back up to last accept
    }
    return out;
}
\end{lstlisting}

The backing up can cause $\Theta(n^2)$ worst-case time: e.g.\ if tokens are \code{a} and
\code{a...ab} the scanner may read far ahead before failing and backing up, repeatedly.
(A memoised version, due to Reps, achieves linear time.)

\section{Simplified maximal munch}

\begin{definition}
\textbf{Simplified maximal munch}: run the DFA until it gets stuck. If the current state
is accepting, emit the token; otherwise report an error. Never back up.
\end{definition}

\begin{lstlisting}[style=cpp]
while (pos < w.size()) {
    int state = START; size_t i = pos;
    while (i < w.size() && delta(state, w[i]) != ERROR)
        state = delta(state, w[i++]);
    if (!accepting(state)) throw ScanError(pos);
    emit(kindOf(state), w.substr(pos, i - pos));
    pos = i;
}
\end{lstlisting}

This runs in $\Theta(n)$ time and is what CS 241's WLP4 and MIPS-assembly scanners use. The
price: some strings that \emph{can} be decomposed into tokens are rejected.

\begin{example}[Simplified maximal munch can fail]
Suppose tokens are $\{a,\ aab\}$ (alternatively: C++ \code{>{}>} vs.\ nested templates). On
input $aaa$: the DFA reads $a$ (accepting), then $aa$ (not accepting, but still alive since
it is a prefix of $aab$), then the third $a$ kills it after $aa$. The state for $aa$ is not
accepting, so simplified munch reports an error. Yet $aaa = a\cdot a\cdot a$ is a valid
decomposition, and full maximal munch would find it by backing up.
\end{example}

\begin{example}[Maximal munch is not always what you want]
In C++03, \code{vector<vector<int>{}>} scanned \code{>{}>} as a single shift token, a famous
annoyance (fixed in C++11 by a special parser rule). Similarly, \code{x-{}-{}-y} is scanned as
\code{x -{}- - y}. Maximal munch is greedy, not a full search over decompositions.
\end{example}

\section{Whitespace, comments, and token separation}

In WLP4, whitespace and comments (\code{//} to end of line) are recognised as ``tokens''
by the DFA and then discarded. Some pairs of tokens must be separated by whitespace,
because under maximal munch they would otherwise merge: \code{int x} vs.\ \code{intx}, or
\code{42abc} (WLP4 treats a NUM immediately followed by a letter as an error). Numbers out of
range (above $2^{31}-1$) are rejected by the scanner as a separate check.

\section{Theory: why not just a regex for the program?}

The language of token sequences is regular, but programming languages require balanced
parentheses, matched braces, and nested structure. $\{(^n )^n\}$ is not regular, so after
scanning we need a more powerful formalism: context-free grammars (next chapter).

\section{Summary}
\begin{itemize}
  \item The scanner turns characters into (kind, lexeme) tokens using a DFA built from the
        regexes for each token kind (union, then subset construction).
  \item Decompositions are ambiguous; maximal munch takes the longest token each time.
  \item General maximal munch backs up to the last accepting state (worst case quadratic).
  \item Simplified maximal munch never backs up: linear time, but may reject decomposable inputs.
  \item Keywords: scan as identifiers then look them up; whitespace and comments are discarded.
\end{itemize}
